% !TeX root = ../../Book3.tex
% 纲要标题：### I.3 射影直线与复数域上的比较
% 状态：新增正文初稿，尚未正式审核。

\section{射影直线与复数域上的比较}
\label{b3:sec:I03}

同一个方程的局部性质与全局覆盖性质并不相同。幂映射在去掉零点和无穷远点后可以平展，但把这两个点补回，便出现分歧。射影直线没有非平凡的连通有限平展覆盖；这条结论在任意代数闭域上成立。复数域还允许用路径追踪覆盖中的点，从环路得到有限纤维上的置换，进而比较拓扑基本群与平展基本群。

\ProseParagraphBreak
这里关于射影直线的证明使用\cref{b3:cor:H-canonical-degree}的典范除子次数公式。复曲线的比较则需要分析理论：拓扑覆盖可以取得复结构，补点后所得有限解析映射还须代数化。后一步采用 Riemann 存在定理的证明概要，并明确列出所用的分析比较结果。

\subsection{代数闭域上的射影直线}
\label{b3:subsec:I03:01}

\begin{proposition}\label{b3:prop:I-projective-line-covers}
	设 \(k\) 为代数闭域。每个非空连通有限平展覆盖
	\[
		f:C\longrightarrow\mathbb P_k^1
	\]
	都是同构。因此，\(\mathbb P_k^1\) 的任意有限平展覆盖都是有限份 \(\mathbb P_k^1\) 的不交并；对任一几何基点 \(\bar x\)，有
	\[
		\pi_1^{\mathrm{\acute et}}(\mathbb P_k^1,\bar x)=1.
	\]
\end{proposition}
\begin{proof}
	先核对源曲线的性质。在 \(\mathbb P_k^1\) 的仿射图上，平展标准模型是 \(k[t,X_1,\ldots,X_r]\) 的局部化模去 \(r\) 条关系，其关于 \(X_1,\ldots,X_r\) 的 Jacobian 行列式可逆。把它视为 \(k\) 代数的展示，便是相对维数一的标准光滑模型。因此 \(C\) 在 \(k\) 上光滑且一维。代数闭域完美，\cref{b3:thm:perfect-field-regular-smooth}使其闭点局部环正则，由\cref{b3:thm:dvr-characterizations}为 DVR。不同不可约分支不能在这样的局部环中相交；有限型性保证分支只有有限个，连通性遂使 \(C\) 为整曲线。

	射影性也需要说明。令 \(L=k(C)\)。有限平展覆盖的局部自由秩是一个正整数，故 \(f\) 满射，且 \(L/k(t)\) 有限。取 \(L\) 的光滑射影曲线模型 \(\bar C\)，这里采用\cite[Chapter I, Theorem 6.9, pp. 44--45]{Hartshorne1977}。函数域嵌入 \(k(t)\hookrightarrow L\) 给出 \(\bar C\) 到 \(\mathbb P_k^1\) 的非恒定有理映射，由\cref{b3:lem:H-rational-map-extension,b3:prop:H-curve-morphism}延拓为有限态射。对 \(\mathbb P_k^1\) 的仿射开集 \(\operatorname{Spec}A\)，原覆盖及这一模型的逆像都是有限正规整 \(A\) 代数的谱，分式域同为 \(L\)。两环都等于 \(A\) 在 \(L\) 中的整闭包：其中元素整于 \(A\)；反之，整于 \(A\) 的元素也整于各该环，由整闭性属于该环。这些识别在局部化下相容，故 \(C\cong\bar C\)。所以 \(C\) 是光滑射影整曲线，可以应用\cref{b3:cor:H-canonical-degree}的次数公式。

	记 \(d=[k(C):k(t)]\)。平展性给出 \(\Omega_{C/\mathbb P^1}=0\)，而\cref{b3:thm:differential-transitivity}在各仿射图上黏合为满射
	\[
		f^*\Omega_{\mathbb P_k^1/k}\longrightarrow\Omega_{C/k}
		\longrightarrow0.
	\]
	两端都是秩一局部自由层；秩一自由模之间的满射由单位相乘给出，因此首箭头为同构。由\cref{b3:exm:H-projective-line-canonical}，
	\[
		\omega_C\cong f^*\omega_{\mathbb P_k^1}
		\cong f^*\mathcal O_{\mathbb P_k^1}(-2).
	\]
	\cref{b3:prop:H-curve-morphism}说明任意点的拉回除子次数为 \(d\)，所以右端线丛的次数为 \(-2d\)。再用\cref{b3:cor:H-canonical-degree}，得到
	\[
		2g(C)-2=-2d,\qquad g(C)=1-d.
	\]
	亏格非负且 \(d\ge1\)，只能有 \(d=1\)。由\cref{b3:prop:H-curve-morphism}的次数一结论，\(f\) 为同构。

	一般有限平展源为 Noether 概形，只有有限多个连通分支；这些分支开闭，限制后仍为非空连通有限平展覆盖。逐个应用已证结论，就得到平凡覆盖的分解。此时纤维函子等同于有限集合的恒等函子，其自然自同构只能为恒等：从单点集合到任意有限集合的各映射迫使每个元素固定。因此平展基本群平凡。
\end{proof}

\cref{b3:exm:I-power-cover}的幂覆盖在补入零点和无穷远点后会出现分歧。对代数闭域 \(k\) 和整数 \(n\ge2\)，当 \(\operatorname{char}k\nmid n\) 时，\(z\mapsto z^n\) 在 \(\mathbb G_{m,k}\) 上是 \(n\) 次有限平展覆盖；延拓到射影直线后，零点和无穷远点的拉回重数都为 \(n\)，所以不再平展。正特征下，即使只去掉一个点，也可能出现新的连通平展覆盖。

\begin{example}\label{b3:exm:I-affine-line-artin-schreier}
	设 \(k\) 为特征 \(p>0\) 的代数闭域。映射
	\[
		\mathbb A_k^1\longrightarrow\mathbb A_k^1,
		\qquad y\longmapsto t=y^p-y
	\]
	对应代数
	\[
		k[t]\longrightarrow k[t,Y]/(Y^p-Y-t)\cong k[y].
	\]
	首一关系使它在 \(k[t]\) 上自由秩为 \(p\)，对 \(Y\) 的导数为 \(-1\)，故由\cref{b3:prop:standard-smooth-lifting,b3:thm:etale-equivalent-characterizations}平展。源为整曲线，所以覆盖连通。平移 \(y\mapsto y+a\)、\(a\in\mathbb F_p\)，给出 \(p\) 个覆盖自同构，并在每个几何纤维上传递，故它是 Galois 覆盖。

	把映射延拓到射影直线，只有无穷远点映到无穷远点。令 \(u=t^{-1}\)、\(v=y^{-1}\)，则
	\[
		f^*u=\frac{v^p}{1-v^{p-1}}.
	\]
	分母在 \(v=0\) 为单位，故该点的分歧指数为 \(p\)。这解释了覆盖为何只在仿射直线上平展；也说明\cref{b3:prop:I-projective-line-covers}中的射影性条件不可去掉。
\end{example}

\subsection{复曲线的有限拓扑覆盖}
\label{b3:subsec:I03:02}

\begin{definition}\label{b3:def:I-topological-cover}
	设 \(M\) 为连通 Riemann 曲面，\(Y\) 为拓扑空间。一个连续映射 \(p:Y\rightarrow M\) 若满足：每个 \(z\in M\) 都有开邻域 \(U\)，其逆像有开集分解
	\[
		p^{-1}(U)=\coprod_{i=1}^n V_i,
	\]
	且每个限制映射 \(p|_{V_i}:V_i\rightarrow U\) 都是同胚，其中 \(n\) 为有限非负整数，就称为\term[youxian-tuopu-fugai]{有限拓扑覆盖}[finite topological covering]。覆盖之间的态射是与到 \(M\) 的映射相容的连续映射。这些对象与态射组成范畴 \(\mathrm{Cov}_{\mathrm{fin}}(M)\)。连通性使纤维基数 \(n\) 与 \(z\) 无关，称为覆盖的次数；\(n=0\) 对应空覆盖。
\end{definition}

连通 Riemann 曲面局部道路连通，因而也道路连通。固定 \(x\in M\)。从 \(x\) 回到 \(x\) 的连续路径，按保持两个端点不变的同伦取等价类，组成\term[tuopu-jibenqun]{拓扑基本群}[topological fundamental group]，记为 \(\pi_1^{\mathrm{top}}(M,x)\)。这里的同伦是连续映射 \(H:[0,1]^2\rightarrow M\)，其上下两边是所比较的路径，左右两边始终取相应端点。固定复合次序：若 \(\alpha,\beta\) 为可拼接的路径，\(\alpha\beta\) 表示先走 \(\beta\)，再走 \(\alpha\)。环路的类以此相乘，常路径为单位，反向路径 \(\bar\alpha\) 给出逆元。拼接的结合性在重新参数化的同伦下成立。
\symidx{\(\pi_1^{\mathrm{top}}(M,x)\)：拓扑基本群，路径乘积先走右因子}

\begin{lemma}\label{b3:lem:I-path-lifting}
	设 \(M\) 为连通 Riemann 曲面，\(p:Y\rightarrow M\) 为有限拓扑覆盖。给定路径 \(\gamma:[0,1]\rightarrow M\) 及 \(y\in p^{-1}(\gamma(0))\)，存在唯一的路径 \(\widetilde\gamma\)，满足 \(p\widetilde\gamma=\gamma\)、\(\widetilde\gamma(0)=y\)。若两条路径在固定端点的条件下同伦，则从同一 \(y\) 出发的提升终点相同。
\end{lemma}
\begin{proof}
	用均匀覆盖邻域覆盖紧集 \(\gamma([0,1])\)，把区间分成有限多个小段，使每段的像落在一个这样的邻域中。起点选定一张覆盖片；在该片中用 \(p\) 的逆同胚提升第一段，再由前一段的终点选定下一张覆盖片。有限次拼接给出提升。任何另一提升都必须在第一小段使用同一张片，随后每段也相同，故唯一。

	对一个路径同伦 \(H:[0,1]^2\rightarrow M\)，把紧正方形细分成有限网格，使每个小方格的像落在一个均匀覆盖邻域内。从给定起点开始，依相邻小方格在相应覆盖片中取逆映射；公共边上的唯一性使这些局部提升相容，得到整个同伦的提升。固定的终点在提升后给出从连通区间到离散纤维的连续映射，只能为常值，因而提升终点不变。
\end{proof}

令 \(\Gamma=\pi_1^{\mathrm{top}}(M,x)\)。对 \(s\in p^{-1}(x)\)，定义 \([\gamma]\cdot s\) 为从 \(s\) 出发提升 \(\gamma\) 的终点。这称为纤维上的\term[danzhi-zuoyong]{单值作用}[monodromy action]。上述引理保证良定义；乘积先走右因子的约定给出
\[
	([\alpha][\beta])\cdot s
	=[\alpha]\cdot([\beta]\cdot s).
\]
所以它是左作用，与\secref{b3:sec:I02}中纤维函子自同构的复合方向一致。

\begin{proposition}\label{b3:prop:I-topological-cover-classification}
	设 \(M\) 为连通 Riemann 曲面，\(x\in M\)，\(\Gamma=\pi_1^{\mathrm{top}}(M,x)\)。取基点纤维并赋予单值左作用，得到范畴等价
	\[
		\mathrm{Cov}_{\mathrm{fin}}(M)
		\simeq \mathrm{FinSet}_{\Gamma}.
	\]
	这里右端为有限左 \(\Gamma\) 集合及等变映射的范畴。非空连通覆盖恰对应非空传递作用。
\end{proposition}
\begin{proof}
	先从一个有限左 \(\Gamma\) 集合 \(S\) 构造覆盖。考虑二元组 \((\gamma,s)\)，其中 \(\gamma\) 从 \(x\) 走到某点 \(z\in M\)，\(s\in S\)。对终点相同的两对，规定
	\[
		(\gamma,s)\sim(\gamma',s')
		\quad\Longleftrightarrow\quad
		s'=[\bar\gamma'\gamma]\cdot s.
	\]
	路径与其反向的连续往返可收缩，故这是等价关系。例如两次变换的作用元相乘后，中间的 \(\gamma'\bar\gamma'\) 消去，得到 \([\bar\gamma''\gamma]\)。令 \(Y_S\) 为这些等价类的集合，投影送到路径终点。

	给定到 \(z\) 的路径 \(\gamma\)，取 \(z\) 的一张小圆盘邻域 \(U\)。对 \(z'\in U\)，用 \(U\) 内从 \(z\) 到 \(z'\) 的路径 \(\delta\) 定义截面
	\[
		z'\longmapsto[(\delta\gamma,s)].
	\]
	圆盘内同端点路径同伦，所以此式与 \(\delta\) 的选择无关。每个投影到 \(U\) 的等价类都唯一落在某个 \(s\) 所给的截面上，故这些截面把逆像分成 \(\#S\) 张片。把各片规定为与 \(U\) 同胚，便定义了 \(Y_S\) 的拓扑；不同圆盘的定义在共同的小圆盘上相容，因此给出有限拓扑覆盖。以常路径识别基点纤维与 \(S\) 后，提升环路的终点正是原来的 \(\Gamma\) 作用。

	反过来，对一个覆盖 \(p:Y\rightarrow M\)，令 \(S=p^{-1}(x)\)。把 \([(\gamma,s)]\) 送到 \(\gamma\) 从 \(s\) 出发的提升终点。\cref{b3:lem:I-path-lifting}和上述等价关系说明此映射良定义。给定 \(y\in Y\)，先把从 \(p(y)\) 到 \(x\) 的一条路径从 \(y\) 提升，再反向行走，就得到从基点纤维到 \(y\) 的路径，故这一映射满射；两个这样的终点相等，反向提升 \(\gamma'\) 就给出 \(s'=[\bar\gamma'\gamma]\cdot s\)，故它也单射。它在每张覆盖片上为同胚，所以是覆盖同构。

	覆盖态射保持路径提升，因而在基点纤维上给出等变映射。每个等变映射 \(u:S\rightarrow T\) 又通过
	\[
		[(\gamma,s)]\longmapsto[(\gamma,u(s))]
	\]
	唯一给出覆盖态射。于是这两个构造在对象和态射上互逆到自然同构，得到所述等价。最后，覆盖空间局部道路连通；基点纤维中两点能在 \(Y\) 中用路径相连，当且仅当它们属于同一个单值作用轨道。每个连通分支都交于基点纤维，所以连通性等价于传递性。
\end{proof}

有限拓扑覆盖还有唯一的相容复结构：在每张覆盖片上拉回 \(M\) 的局部复坐标即可。两张图的过渡函数就是底曲面的全纯过渡函数，因此这一构造确实给出 Riemann 曲面，并使 \(p\) 局部双全纯。覆盖态射在这些图中也自动全纯。有限张覆盖片还使 \(p\) 为固有映射：给定底上的紧集，用有限个更小圆盘覆盖它，使每个小圆盘的闭包落在一个均匀覆盖圆盘内。各闭包逆像是有限张紧片，紧集的逆像为这些紧片之并中的闭子集，故仍紧。于是有限拓扑覆盖与有限平展解析覆盖承载同一组对象和态射。

\subsection{Riemann 存在定理与基本群比较}
\label{b3:subsec:I03:03}

对复代数曲线 \(X\)，把仿射图中的多项式方程视为复解析方程，再按原过渡映射黏合，得到它的\term[jiexihua]{解析化}[analytification] \(X^{\mathrm{an}}\)。光滑性使它成为 Riemann 曲面；连通代数曲线的解析化也连通，这里采用代数与解析连通分支的比较\cite[Exposé XII, Corollaire 2.6, p. 244]{SGA12003}。

\begin{theorem}[Riemann 存在定理]\label{b3:thm:I-riemann-existence}
	设 \(X\) 为光滑连通、有限型、分离的复代数曲线。解析化给出范畴等价
	\[
		\operatorname{FinEt}(X)
		\xrightarrow{\ \sim\ }
		\mathrm{Cov}_{\mathrm{fin}}(X^{\mathrm{an}}).
	\]
	它保持基点纤维；这里允许 \(X\) 不是射影曲线。这一比较称为\term[Riemann-cunzai-dingli]{Riemann 存在定理}[Riemann existence theorem]。
\end{theorem}
\begin{proofsketch}
	有限平展代数态射的解析化仍有限平展；局部平展标准模型的可逆 Jacobian 与复逆函数定理说明，它在解析拓扑中局部双全纯，故为有限拓扑覆盖。反向需要把解析覆盖代数化。以下给出曲线情形的构造；一般的比较定理见\cite[Exposé XII, Théorème 5.1, pp. 251--253]{SGA12003}。

	取光滑射影紧化 \(X\subseteq\bar X\)，其补集 \(D\) 为有限个点。这样的紧化来自曲线的光滑射影模型，见\cite[Chapter I, Corollary 6.10, p. 45]{Hartshorne1977}。给定有限拓扑覆盖 \(Y\rightarrow X^{\mathrm{an}}\)，先赋予上一小节的唯一复结构。对每个补点，取局部坐标圆盘 \(\Delta\)，使它不含其他补点。径向变形把去心圆盘 \(\Delta^*\) 缩到圆周；连续选择辐角后，一条闭环路的辐角总变化是 \(2\pi\) 的整数倍。该整数为零时，辐角函数的直线同伦使环路收缩；每个整数又都可由相应圈数的环路实现。因此 \(\pi_1^{\mathrm{top}}(\Delta^*)\cong\mathbb Z\)，沿小圆正向绕行一次给出生成元。有限纤维在该生成元作用下分成循环轨道，每个长度 \(e\) 的轨道对应连通覆盖
	\[
		\Delta^*\longrightarrow\Delta^*,\qquad w\longmapsto z=w^e.
	\]
	这先由拓扑覆盖分类得到；覆盖同构在拉回复结构下自动全纯，所以也是解析同构。对每个循环补入 \(w=0\)，覆盖便延拓成有限全纯映射
	\[
		\bar Y\longrightarrow\bar X^{\mathrm{an}},
	\]
	其中 \(\bar Y\) 是有限个紧 Riemann 曲面的不交并。局部指数 \(e\) 记录新补点处的分歧，原来开曲线上的映射仍无分歧。

	紧化后的底曲线是固有代数概形。此时采用\cite[Exposé XII, Corollaire 4.6, p. 251]{SGA12003}的有限映射比较，把 \(\bar Y\) 连同到 \(\bar X^{\mathrm{an}}\) 的有限解析映射代数化。这个结果允许补点处有分歧：有限映射等同于底上的凝聚代数，固有情形的凝聚层比较先代数化其底模，再由完全忠实性代数化乘法与单位。把所得有限代数映射限制到 \(X\)，由代数与解析平展性的比较\cite[Exposé XII, Proposition 3.1(iii), pp. 244--245]{SGA12003}得到有限平展覆盖。这证明本质满性。

	对完全忠实性，考虑两个覆盖间的拓扑态射。它已经全纯；在补点旁，它从指数 \(e_1\) 的循环轨道映到指数 \(e_2\) 的循环轨道，等变性使 \(e_2\mid e_1\)。选定局部坐标后，映射具有
	\[
		w_2=\zeta w_1^{e_1/e_2},\qquad \zeta^{e_2}=1
	\]
	的形式，因此唯一延拓到所补的点。固有底上的有限映射比较把这一延拓唯一代数化，限制后得到原来开曲线上的代数态射。由此得到态射集合的双射；基点纤维在各步均是原来的有限点集。

	这里采用了复逆函数定理、去心圆盘的解析覆盖模型，以及固有情形的凝聚层比较和有限映射代数化。后两个比较定理涉及的分析层与代数层理论不在本书展开；上述构造说明它们在存在性、态射延拓和平展性核验中的具体作用。
\end{proofsketch}

若 \(x\in X(\mathbb C)\)，它确定几何基点 \(\bar x:\operatorname{Spec}\mathbb C\rightarrow X\)。上面的等价与取纤维相容，表现为
\[
	\begin{tikzcd}[column sep=large]
		\operatorname{FinEt}(X)
		\arrow[r,"\operatorname{an}","\sim"']
		\arrow[dr,"F_{\bar x}"']
		&
		\mathrm{Cov}_{\mathrm{fin}}(X^{\mathrm{an}})
		\arrow[d,"F_x"]\\
		&\mathrm{FinSet}.
	\end{tikzcd}
\]
因此，每条以 \(x\) 为基点的环路在全部代数有限平展覆盖的纤维上给出相容置换。由路径乘积约定，这些置换定义群同态
\[
	\rho:\pi_1^{\mathrm{top}}(X^{\mathrm{an}},x)
	\longrightarrow\pi_1^{\mathrm{\acute et}}(X,\bar x).
\]

\begin{corollary}\label{b3:cor:I-comparison-fundamental-groups}
	设 \(X\) 为光滑连通、有限型、分离的复代数曲线，\(x\in X(\mathbb C)\)，\(\bar x\) 为相应几何基点。记 \(\Gamma=\pi_1^{\mathrm{top}}(X^{\mathrm{an}},x)\)。则 \(\rho\) 诱导拓扑群的典范同构
	\[
		\widehat\Gamma
		\cong\pi_1^{\mathrm{\acute et}}(X,\bar x),
		\qquad
		\widehat\Gamma
		=\varprojlim_{\substack{N\triangleleft\Gamma\\{}[\Gamma:N]<\infty}}\Gamma/N.
	\]
	左端称为 \(\Gamma\) 的\term[toushe-youxian-wanbeihua]{投射有限完备化}[profinite completion]，赋予有限离散商群逆极限的拓扑。
\end{corollary}
\begin{proof}
	有限 \(\Gamma\) 作用的核是正规有限指数子群，所以每个有限左 \(\Gamma\) 集合的作用都唯一延拓为连续左 \(\widehat\Gamma\) 作用。唯一性来自 \(\Gamma\) 的自然像稠密：有限多个商群中的相容坐标可在它们对应子群的交所给的单个有限商中表示，再取一个 \(\Gamma\) 中的代表。反之，有限离散集合上的连续 \(\widehat\Gamma\) 作用限制后就是有限 \(\Gamma\) 作用。等变映射也相同，因为可以在两个作用核的交所给的有限商上核验。因此这两个有限作用范畴连同遗忘函子相同。

	还可直接恢复遗忘函子的自同构群。对每个正规有限指数子群 \(N\)，考虑左正则集合 \(\Gamma/N\)。自然自同构必须与其全部右平移相容，故在这个集合上是某个元素 \(g_N\in\Gamma/N\) 的左平移；自然性使 \((g_N)_N\) 在各商映射下相容。每个有限 \(\Gamma\) 集合都是若干个这样的正则集合的等变像：先取其作用核 \(N\)，再对每个元素 \(s\) 使用 \(gN\mapsto g\cdot s\)。所以这些左平移唯一确定整个自然自同构，反过来每个相容族也确实作用于全部有限 \(\Gamma\) 集合。由此得到遗忘函子的自同构群为 \(\widehat\Gamma\)，逐点稳定有限纤维的拓扑正是上述逆极限拓扑。

	\cref{b3:prop:I-topological-cover-classification,b3:thm:I-riemann-existence}把这一遗忘函子识别为 \(F_{\bar x}\)。而\cref{b3:thm:I-galois-classification}中的平展基本群就是 \(\operatorname{Aut}(F_{\bar x})\)，故得到所述拓扑同构；对 \(\Gamma\) 中的元素，它由环路提升给出，因此延拓的是 \(\rho\)。
\end{proof}
\symidx{\(\widehat\Gamma\)：离散群 \(\Gamma\) 的投射有限完备化}

这一比较的来源是\cite[Exposé XII, Corollaire 5.2, pp. 253--254]{SGA12003}。有限覆盖只能检测基本群的有限商；上式中的完备化正是同时保留这些有限商，而不把离散拓扑基本群直接当作平展基本群。

\begin{example}\label{b3:exm:I-complex-curve-groups}
	射影直线的复点组成球面 \(S^2\)，其拓扑基本群平凡。可将一条环路在有限张小圆盘中同伦为分段光滑路径，再选一个不在这有限条弧上的点；去掉该点得到平面，环路遂可收缩。由\cref{b3:cor:I-comparison-fundamental-groups}，
	\[
		\pi_1^{\mathrm{\acute et}}(\mathbb P_{\mathbb C}^1,\bar x)=1,
	\]
	与\cref{b3:prop:I-projective-line-covers}的任意特征代数证明一致。

	对 \(\mathbb G_{m,\mathbb C}^{\mathrm{an}}=\mathbb C^*\)，径向变形把它收缩到单位圆，环路的绕数给出
	\[
		\pi_1^{\mathrm{top}}(\mathbb C^*,1)\cong\mathbb Z,
		\qquad
		\pi_1^{\mathrm{\acute et}}(\mathbb G_{m,\mathbb C},\bar1)
		\cong\widehat{\vphantom{Z}\mathbb Z}.
	\]
	具体地，环路 \(t\mapsto e^{2\pi it}\) 在 \(z\mapsto z^n\) 下，从 \(e^{2\pi ij/n}\) 出发的提升为 \(t\mapsto e^{2\pi i(j+t)/n}\)。终点把 \(j\) 送到 \(j+1\)，所以这份覆盖对应 \(\mathbb Z/n\mathbb Z\) 上的循环左作用。每个非空传递有限 \(\mathbb Z\) 集合都是一个这样的循环，故所有连通有限平展覆盖都同构于某个 \(n\) 次幂覆盖。

	再取 \(U=\mathbb P_{\mathbb C}^1\setminus\{0,1,\infty\}\)。在平面上，先用大圆围住 \(0,1\)，再用两个小圆分别围住它们；去心内圆盘和大圆外部可沿径向缩回边界，所得三边界曲面又可缩回两个圆加一条连接弧的图。取连接弧为生成树，便得到两个自由环路，所以
	\[
		\pi_1^{\mathrm{top}}(U^{\mathrm{an}},x)\cong F_2,
		\qquad
		\pi_1^{\mathrm{\acute et}}(U,\bar x)\cong\widehat{F_2}.
	\]
	因此其 \(n\) 次连通覆盖可由一对在 \(n\) 个字母上生成传递作用的置换描述，更换纤维标号对应同时共轭。取两个不同的对换就得到三点集合上的传递 \(S_3\) 作用；对应的三次覆盖不是 Galois 覆盖，而其 Galois 闭包为六次覆盖。这与 \(\mathbb C^*\) 的循环覆盖已经有不同的有限群结构。
\end{example}


% BEGIN APPENDIX OUTLINE SYNC
% * 状态：初稿已完成；新增范围尚未正式审核。
% #### I.3.1 代数闭域上的射影直线
% #### I.3.2 复曲线的有限拓扑覆盖
% #### I.3.3 Riemann 存在定理与基本群比较
% END APPENDIX OUTLINE SYNC
% ---
